\chapter[Conclusion]{Conclusion}
\label{cp:conclusion}

This thesis set out to show that the dependency relationships recorded in a Software Bill of Materials are an analytical resource, not just a record. The SBOM tools surveyed in \autoref{sec:sbom-tools} treat an SBOM as an inventory, at most tracing the paths to a component or listing its dependents; SBOM Lens builds a directed dependency graph and computes dominance and structural measures over it. This chapter summarises what the evaluation found (\autoref{sec:concl-findings}), revisits the research objectives of \autoref{sec:objectives} against the contributions made (\autoref{sec:concl-contributions}), states the limitations of the work (\autoref{sec:concl-limitations}), and identifies directions for future work (\autoref{sec:concl-future}).


\section{Summary of Findings}
\label{sec:concl-findings}

The evaluation of \autoref{cp:evaluation} applied SBOM Lens to the production dependency graphs of 90 popular npm applications, each with at least one critical known vulnerability, and answered three research questions.

\textbf{RQ1: the dependency graph of a single application has the heavy-tailed structure found in whole ecosystems.} The median graph has 332 components, of which only 11\% are declared by the application itself. A power law could be fitted to the in-degree distribution of 79 of the 90 graphs, with exponents between 2.09 and 2.99 (median 2.58), and every graph is small-world against a weak random baseline: components form tightly clustered neighbourhoods while remaining about five undirected hops apart. Version conflicts are pervasive (median 24 per application), and 30 applications contain dependency cycles. The structural measures of \autoref{cp:network-sci} therefore have structure to work with at the scale of one application.

\textbf{RQ2: severity says almost nothing about where a vulnerable component sits.} Ranking each application's vulnerable components by severity band and by structural position gives essentially unrelated orderings: the median Kendall \(\tau_b\) is 0.00 both for blast-radius size and for HITS authority (in the 79 graphs where it is defined), so that in most applications none of the critical components is among the equally many components with the largest blast radius. A critical vulnerability is as likely to affect a package with a single consumer as a component through which much of the application is reached.

\textbf{RQ3: most vulnerability exposure is structural, and a flat inventory does not show how it arises.} In the median application, 75\% of the vulnerable components are not declared by the application; every application has vulnerable chokepoints --- vulnerable components that parts of the application can be reached only through --- and in the median application they are 45\% of the vulnerable components and dominate 21\% of the graph; and 19\% of the application's components are vulnerable or depend on a vulnerable component (12\% depend on one).

Together, these results support the central claim of the thesis: the position of a component in the dependency graph carries risk information that a flat inventory of components and vulnerabilities does not, and SBOM Lens makes that information available at interactive speed: on graphs of up to a thousand components, every analysis except the audit took milliseconds in the median case and under three seconds at most; the audit, which also queries the npm registry once per package, took up to 16 seconds.


\section{Contributions}
\label{sec:concl-contributions}

The four research objectives of \autoref{sec:objectives} map one-to-one to the contributions of this work; each was met, with the qualification stated for RO3.

\textbf{Graph construction pipeline (RO1).} SBOM Lens ingests CycloneDX JSON SBOMs and constructs a directed dependency graph in which every component is a node and every declared dependency a consumer-to-dependency edge (\autoref{cp:system-design}). The graph is built once at ingestion, cached, and shared by every analysis, and it is rendered interactively with Cytoscape.js, where the analyst can explore, edit and export it. Edits are written back to the backend, which rebuilds the graph and invalidates every cached analysis, so the structural results always describe the current graph.

\textbf{Vulnerability integration (RO2).} Every component with a Package URL is matched against OSV, which aggregates the GitHub Advisory Database and other open-source advisory sources (\autoref{cp:vuln-analysis}). Severity is derived from the CVSS v3 vector by a client-side implementation of the v3.1 specification, falling back, when no such vector is present, to the advisory's own severity label, and otherwise recorded as unknown rather than silently treated as low. Vulnerable components are coloured on the graph by their highest severity band, and for any vulnerable component the analyst obtains its blast radius, up to five shortest dependency paths through which it enters the application, and its dominated set --- the components that could not be reached without it.

\textbf{Graph-based structural risk analysis (RO3).} SBOM Lens computes structural measures at three scales (\autoref{cp:network-sci}): per component (articulation points and bridges, HITS scores, a spectral exposure score, core numbers), per cluster (Louvain communities, the bow-tie decomposition, the rich-club coefficient), and for the whole graph (the spectral radius and epidemic threshold, degree entropy, cycles, version conflicts, namespace concentration, the in-degree distribution and the small-world coefficient). For each measure, the thesis states precisely what is computed, on which graph, and how the result should be read --- including where a measure is only a coarse indicator. The evaluation shows that these measures carry information that severity does not; it does not show that prioritising by them leads to better remediation, which would need exploitation data or analyst outcomes.

\textbf{Empirical evaluation (RO4).} The evaluation of \autoref{cp:evaluation} applies SBOM Lens to the production dependency graphs of npm applications, collected reproducibly from public repositories at fixed release tags and restricted to applications with at least one critical known vulnerability. The collection pipeline, the SBOMs, every analysis result and the analysis script are published in the \texttt{sbom-lens-eval} repository, from which the tables and figures of \autoref{cp:evaluation} can be regenerated.
On 90 applications, the structural and severity rankings of vulnerable components were found to be essentially unrelated (median \(\tau_b = 0.00\)), and three-quarters of the vulnerable components to be transitive-only.

Beyond the objectives, the work contributes a clarification of the vocabulary of structural SBOM analysis: it distinguishes the blast radius of a component from its dominated set, chokepoints from articulation points, version conflicts from the diamond pattern that causes them, and static spectral analysis from propagation simulation. Where the interface uses a name that suggests more than is computed --- the \texttt{r0} field, the \emph{Diamond Dependencies} label --- the thesis states what is actually computed.


\section{Limitations}
\label{sec:concl-limitations}

\textbf{Vulnerability data.} SBOM Lens relies on OSV as its only vulnerability source. OSV aggregates the main open-source advisory databases, but an advisory reaches SBOM Lens only once OSV has ingested it, and vulnerabilities recorded only in NVD, without an open-source advisory, are not reported. Severity is computed from CVSS v3 vectors only; records that carry only a CVSS 4.0 vector fall back to the advisory's own label or to the unknown band. Exploit likelihood (EPSS) and known-exploitation catalogues are not integrated, so severity describes potential impact rather than the probability of exploitation.

\textbf{Embedded vulnerability data (VEX).} SBOM Lens ignores the \texttt{vulnerabilities} array of a CycloneDX document (\autoref{sec:cyclonedx-structure}), which may carry Vulnerability Exploitability eXchange statements from the software's producer. Querying OSV directly gives every analysis the same vulnerability data, current at upload, regardless of what the producer chose to embed, and avoids trusting producer-supplied assessments. The cost is that a producer's statement that a vulnerability is not exploitable in their product --- often the most useful information for prioritisation --- is not taken into account, so SBOM Lens may report vulnerabilities that the producer has already assessed as irrelevant.

\textbf{Component-level analysis.} All analysis is at the level of components. A component in the blast radius of a vulnerability depends on the vulnerable component, but may never execute the vulnerable code (\autoref{sec:rw-call-graphs}); exposure is therefore over-approximated.

\textbf{Directed and undirected measures.} Several measures --- articulation points, bridges, core numbers, communities, the rich-club coefficient, the small-world coefficient and all spectral measures --- are computed on the undirected projection of the dependency graph, either because they are defined only for undirected graphs or because, as for the spectral radius, they vanish on acyclic directed graphs (\autoref{cp:network-sci}). These measures describe how densely components are interconnected, not the direction in which dependencies run, and they must be read accordingly.

\textbf{Heuristic thresholds and baselines.} Some outputs rest on conventions rather than theory: the labels of the degree entropy, the 60\% namespace-concentration flag, the 0.5 rich-club threshold, and the fan-in threshold for abstraction candidates. The small-world coefficient uses a closed-form random baseline that ignores the degree sequence, the rich-club coefficient is not normalised against a null model, and the power-law fit is not compared against alternative distributions. These outputs are reported as descriptive indicators, not as statistical tests.

\textbf{Numerical edge cases.} When HITS power iteration does not converge, SBOM Lens reports every score as zero, and when no valid power-law fit exists it returns the package's bound \(\alpha = 3\); neither case is flagged. In the evaluation each occurred in 8 graphs, and the HITS results, which rest on 79 of the 90 graphs, depend on the NetworkX version (\autoref{sec:hits}).

\textbf{Evaluation sample.} The evaluation covers popular npm applications only, not organised as monorepos, at release tags on or before 31 March 2024, with a committed lockfile, between 50 and 1{,}000 components in their production dependency graph, and at least one critical known vulnerability. Each of these filters shapes the sample, and so does the SBOM generator, which produced no usable graph for all but one \texttt{pnpm} project and an incomplete one, whose counts are lower bounds, for seven applications. The results describe vulnerable, medium-sized npm applications and may not transfer to other ecosystems, where duplicate versions are rarer, or to very large applications. Vulnerability data reflects OSV on 30 September 2026, so it includes advisories published after each release, which an analyst could not have seen at the time: 93\% of the findings, and every critical finding in 60 of the 90 applications (\autoref{sec:eval-threats}).

\textbf{Implementation.} SBOM Lens is a research prototype. It keeps all state in memory, so SBOMs and vulnerability records are lost when the backend restarts; it has no authentication and is intended as a single-user local tool; it accepts CycloneDX JSON only; and it queries OSV only when an SBOM is uploaded, so vulnerability records are not refreshed and a component added by editing is not looked up.


\section{Future Work}
\label{sec:concl-future}

\textbf{Propagation simulation.} SBOM Lens analyses the dependency graph statically. A natural extension is to simulate how a compromise spreads over time, with an epidemic model such as SIR run on the dependency graph (\autoref{sec:propagation-modelling}). The spectral radius already computed by SBOM Lens would then serve as a check on the simulation, and the simulated reach could be compared with the static vulnerability reach reported in \autoref{cp:evaluation}.

\textbf{Richer vulnerability data.} Integrating EPSS scores and known-exploitation catalogues would let structural measures be combined with the likelihood of exploitation, not only with severity. Direct integration of NVD, despite the imprecision of CPE matching (\autoref{sec:multi-source-rationale}), and of the GitHub Advisory Database by version range would reduce dependence on OSV's aggregation, and a CVSS 4.0 calculator would remove the unknown band for records that carry only a version-4 vector. Consuming VEX statements from the SBOM would let the producer's exploitability assessments suppress irrelevant findings.

\textbf{Reachability.} Combining the component-level graph with call-graph reachability (\autoref{sec:rw-call-graphs}) would prune the vulnerability overlay to vulnerabilities whose code the application can actually execute, where source code is available.

\textbf{Broader evaluation.} Replicating the evaluation for PyPI and Maven applications would test whether the findings hold in ecosystems with different dependency resolution rules. A user study with security analysts would test the question the metric-based evaluation cannot: whether the structural view changes, and improves, the decisions analysts actually make.

\textbf{Platform.} Persistent storage, support for SPDX input, and a version history for each application, extending the pairwise SBOM comparison view outside the scope of this thesis (\autoref{sec:architecture-overview}), would turn the prototype into a tool for continuous use, in which the structural risk of an application could be tracked as its dependencies evolve.
